% Abstract (English, một trang): bối cảnh, mục tiêu, phương pháp, kết quả chính,
% giới hạn. Không dùng hình, bảng hay trích dẫn.
\FrontChapter{ABSTRACT}

A single amino acid substitution can alter the activity of a degron, a sequence element involved in
selective protein degradation. This thesis predicts substitution effects as changes in the Protein
Stability Index ($\Delta$PSI) measured by a peptide reporter assay and develops a local analysis tool.
The computational work reconstructs labels, controls sequence overlap and evaluates models with
protein-group-aware folds, inner model selection and group-bootstrap confidence intervals.

The full-PSI development set contains \nVthreeNPeptides{} peptides, \nVthreeNCells{} mutation cells
and \nVthreeNGroups{} protein groups. Expanding ridge training from WT PSI at most 4 to the full
range reduces peptide-macro mean absolute error from \nVthreeRidgeB{} to \nVthreeRidgeD{} PSI.
Frozen ESM2-35M with a convolutional head reaches \nVthreeCnnD{} MAE; its matched pointwise
control reaches \nVthreePointwiseD{}. Scanning-label pre-training improves a one-hot CNN under
the tested design, whereas quality weighting worsens ridge. The lowest selected MAE does not
establish an architecture winner. Earlier restricted-domain experiments also provide a calibrated
ridge baseline and numerical agreement between its inference and research paths.

As of 3 November 2026, a React/FastAPI baseline exists, while persistent projects, batch/protein
jobs and experimental-data comparison remain incomplete. The newer workspace is a synthetic-data
review prototype. The local pLM and LoRA expansion, full-PSI calibration and locked held-source
evaluation are unfinished; the serving model has not been replaced. All predictive findings are
development evidence from the assay, with no claim of independent external validation, intact-protein
effects, E3-specific recognition or clinical interpretation.

\par\vspace{6bp}\noindent\textbf{Keywords:} degron, protein stability, mutation effect prediction, ubiquitin--proteasome system, machine learning.
